% Quantitative consequences of the strict seven-cycle interval certificate.
\subsection{Parameter-uniform isolation and stability}\label{sec:consequences}
Write $\theta=(A,B,\mu)$ and
$\theta_0=(9,3/2,181363/100000)$.
Let $x_\theta(t;y)$ solve the initial-value problem with
$x_\theta(0;y)=y$, and put
$D_\theta(y)=x_\theta(2\pi;y)-y$.

Theorem~\ref{thm:main} already states the uniform result.
Its proof reduces to three checks: preservation of the
six Poincar\'e displacement signs, absence of grazing on
the three shooting intervals, and the hyperbolicity of
the central antisymmetric orbit.

\begin{proof}[Completion of the proof of Theorem~\ref{thm:main}]
We use the three ingredients described above, with all
inequalities uniform over the cube~\eqref{eq:param-cube}.

\emph{1. Uniform displacement control.}
The reference vector field points strictly inward on the
boundaries $x=\pm2$, since $15/2-\mu_0>0$.
Consequently $|x_{\theta_0}(t;y)|\le2$ for all
$|y|\le2$ and $t\ge0$.
For $\theta$ in \eqref{eq:param-cube}, all secant slopes
of the perturbed scalar field $f_\theta(\cdot,t)$ are
at most $B<2$, whereas on each reference orbit
\[
 |f_\theta(x_{\theta_0},t)-f_{\theta_0}(x_{\theta_0},t)|
 \le|\Delta A|(|x_{\theta_0}|+|\sat(x_{\theta_0})|)
 +|\Delta B||\sat(x_{\theta_0})|+|\Delta\mu|
 <5\cdot10^{-12}.
\]
The scalar one-sided Gronwall inequality for solutions
started from the same initial value thus gives
\begin{equation}\label{eq:uniform-parameter}
 |x_\theta(t;y)-x_{\theta_0}(t;y)|
 \le5\cdot10^{-12}\frac{e^{2t}-1}{2}
 <1.5\cdot10^{-6},\quad
 0\le t\le2\pi,\quad |y|\le2.
\end{equation}
Indeed, $\pi<13/4$ and $e<11/4$ imply
$e^{4\pi}<e^{13}<(11/4)^{13}<600000$.
The smallest absolute certified displacement endpoint
margin in Table~\ref{tab:signs} exceeds
$4.96788\cdot10^{-6}$, so all six signs persist.
The intermediate value theorem therefore gives at least
one periodic solution in each of $I_1,I_2,I_3$.
At the central parameter triple, their existence and
the strict switching pattern are established by
Proposition~\ref{prop:witness}.

\emph{2. Explicit grazing barriers and isolation.}
At a tangency to $\ell\in\{-1,1\}$ the phase
$t_*$ must satisfy $\mu\sin t_*+B\ell=0$.
For the parameters under consideration $\mu>B>0$.
There are exactly four such phases in a period,
and each lies within $10^{-10}$ of its reference
phase: the perturbation of the phase equation
at the old root is less than $2\cdot10^{-12}$,
whereas the modulus of its time derivative
is greater than $0.8$ in the relevant
$10^{-10}$ neighborhood.

The original exact-flow directed-interval calculation
for the four phases and the \emph{six} rational shooting
endpoints, as now checked by
\texttt{proof/certify\_parameter\_persistence.py},
establishes the stronger quantitative bound
\begin{equation}\label{eq:grazingmargin}
 \bigl|x_{\theta_0}(t_*^{\,0};y)-\ell\bigr|
 >3\cdot10^{-4}
\end{equation}
at every relevant phase $t_*^{\,0}$, endpoint $y$,
and associated threshold $\ell$.
This inequality concerns absolute distances; their
signs are those in the grazing-barrier table of
Section~\ref{sec:proof}.
The scalar velocity is bounded by $40$ on the
trajectories under consideration, since
$|x_\theta|<2.000002$. Combining this with
\eqref{eq:uniform-parameter} gives
\[
 \bigl|x_\theta(t_*^\theta;y)-
              x_{\theta_0}(t_*^{\,0};y)\bigr|
 <1.5\cdot10^{-6}+40\cdot10^{-10}
 <1.505\cdot10^{-6}.
\]
Thus every strict barrier sign persists at
each endpoint. Since $y\mapsto x_\theta(t;y)$
is strictly increasing, none of the trajectories
with $y\in I_j$ can graze a switching level.
The shooting trajectories also stay uniformly
away from switching levels at $t=0,2\pi$;
there can be no crossing loss at the temporal
endpoints. Continuing parameters inside the
connected cube preserves the transverse
crossing itinerary on each $I_j$. Hence
$D_\theta$ is real analytic in $y$ on a
neighborhood of every $I_j$.
Opposite endpoint signs prevent the analytic
function from being identically zero.
Its zeros on $I_j$ are finite and isolated.

The three endpoint sign patterns are
$(+,-),(-,+),(+,-)$, respectively.
A positive-to-negative zero of a scalar
return displacement is locally attracting,
and a negative-to-positive zero is locally
repelling: the strictly increasing Poincar\'e
map moves nearby iterates monotonically
toward, respectively away from, the fixed point.
Therefore at least two attracting and
one repelling isolated periodic solutions
lie above the central antisymmetric solution.

\emph{3. Quantitative persistence of the central repeller.}
For all $\theta$ the equation has a unique
antisymmetric periodic orbit. Its initial
value $y_c(\theta)$ is the unique root of
$Q_\theta(y)=x_\theta(\pi;y)+y$, a strictly
increasing function with incremental slope
at least one. The base certificate encloses
$y_c(\theta_0)$ between
$-0.579022471$ and $-0.579022470$.
Using the improved bound $B<1.51$ up to
$t=\pi$, we have
\[
 |x_\theta(\pi;y)-x_{\theta_0}(\pi;y)|
 \le5\cdot10^{-12}\frac{e^{1.51\pi}-1}{1.51}
 <5\cdot10^{-12}\frac{158-1}{1.5}
 <6\cdot10^{-10},
\]
where $1.51\pi<5$ and $e^5<(11/4)^5<158$.
Consequently
$|y_c(\theta)-y_c(\theta_0)|<6\cdot10^{-10}$.
For $0\le t\le6/5$, the same estimate and
dependence on initial data yield
\[
 |x_\theta(t;y_c(\theta))-
        x_{\theta_0}(t;y_c(\theta_0))|
 < 5\cdot10^{-9},
\]
using $e^{1.51(6/5)}<e^2<(11/4)^2$.

The strict reference center shooting brackets,
together with monotonicity of downward and
upward crossing times in the shooting initial
value, place the two lower-zone events in
\[
 0.63458137<t_1^0<0.63458138,\qquad
 1.08317132<t_2^0<1.08317133.
\]
The threshold velocities $\mu_0\sin t-B_0$
are less than $-0.4$ on the first box and
greater than $0.09$ on the second.
For $|t-t_j^0|\le10^{-6}$, the absolute
change in $\dot x$ is less than
$410\cdot10^{-6}<0.001$, because
$|\dot x|<40$, $|f_x|<10$, and $\mu<2$.
Hence the reference level gaps at
$t_j^0\pm10^{-6}$ have absolute value
at least $0.089\cdot10^{-6}>5\cdot10^{-9}$.
Both crossing signs persist for the perturbed
central orbit, and its two crossing times
differ from the reference ones by less
than $10^{-6}$. The absence of further
lower-zone crossings follows from the
single-positive/single-negative phase arcs
of the threshold velocity (Lemma~\ref{lem:cross});
the upper level cannot be crossed and
recrossed on $[0,\pi]$, where its vector field
points upward. Thus its lower-zone residence
$s_\theta$ in the half-period obeys
\[
 s_\theta<(1.08317133-0.63458137)+2\cdot10^{-6}
 =0.44859196<0.448592.
\]
By antisymmetry the total exterior residence
in one period is $2s_\theta$, and the exact
Floquet exponent is
\begin{align*}
 \log P_\theta'(y_c(\theta))
 &=2\pi B-2(A+B)s_\theta\\
 &>2(3.14159)(1.5-10^{-12})
     -2(10.5+2\cdot10^{-12})(0.448592)
 >0.0043>0.
\end{align*}
The central orbit is therefore hyperbolic
repelling and isolated throughout the cube.

Finally, half-period reflection
$x(t)\mapsto-x(t+\pi)$ preserves each local
stability type, and provides three distinct
isolated counterparts below the central
orbit. This gives at least four attracting
and three repelling periodic orbits.
\end{proof}

\begin{corollary}[Critical residence-time profiles forced by the cycles]
\label{cor:critical}
For every parameter triple in \eqref{eq:param-cube},
the mean-parametrized biased response has at least
three distinct positive-mean stationary profiles
satisfying
\[
   \Lambda_\xi(\xi,\mu)=0,\qquad
   T(\xi,\mu)=\rho=\frac{2\pi A}{A+B}.
\]
\end{corollary}
\begin{proof}
Theorem~\ref{thm:main} gives three ordered
periodic orbits above the antisymmetric central
solution, with means $0<\xi_1<\xi_2<\xi_3$.
The central mean is zero and the biased
response vanishes at all four means.
Rolle's theorem applied to $(0,\xi_1)$,
$(\xi_1,\xi_2)$, and $(\xi_2,\xi_3)$
gives three distinct stationary profiles.
The identity \eqref{eq:key}, including the
continuous extension at an isolated grazing,
gives $T=\rho$ at all three.
Neither nondegeneracy of these points nor
an exact global count of folds is asserted.
\end{proof}

The supplemental script
\texttt{proof/certify\_parameter\_persistence.py}
verifies the directed interval margins in
\eqref{eq:grazingmargin}, the central event brackets,
and the rational inequalities used here. It reuses
the original directed-\texttt{Decimal} arithmetic;
it is not an independent Arb/FLINT implementation.
